\subsection{Partitions of the spectrum of an element}

Let A be a unital Banach algebra, \(x \in A\) , and \(K = \operatorname{Sp}_{\mathrm{A}}(x)\) . Let \(\Pi\) denote the set of subsets of K that are both open and closed in K. Let B be the closed full subalgebra of A generated by x; it is commutative.

For every \(H \in \Pi\) , there exists a unique element \(f_{H}\) of \(\mathcal{O}(K)\) equal to 1 in a neighborhood of H and to 0 in a neighborhood of K---H. We set \(j_{\mathrm{H}} = f_{\mathrm{H}}(x)\) . The element \(j_{H}\) is an idempotent of A, called associated with x and H, and we have the following formulas:

\[
j _ {\mathrm{H} \cap \mathrm{H} ^ {\prime}} = j _ {\mathrm{H}} j _ {\mathrm{H} ^ {\prime}} = j _ {\mathrm{H} ^ {\prime}} j _ {\mathrm{H}}\tag{15}
\]

(H, H' ∈ Π)

\[
j _ {\mathrm {H\cup H^ {\prime}}} = j _ {\mathrm{H}} + j _ {\mathrm {H^ {\prime}}} - j _ {\mathrm {H^ {\prime}}} j _ {\mathrm{H}}\tag{16}
\]

\[
(\mathrm{H}, \mathrm{H} ^ {\prime} \in \Pi)
\]

\[
j _ {\emptyset} = 0, \qquad j _ {\mathrm{K}} = 1.
\]

Let \(H \in \Pi\). We define \(A_{H} = j_{H}Aj_{H}\). This is a closed subalgebra of A, admitting the identity element \(j_{H}\) (cf.~Lemma 1 of I, p.~2). One also sets \(B_{H} = j_{H}Bj_{H}\) and \(x_{H} = xj_{H} = j_{H}x = j_{H}xj_{H} \in B_{H}\) .

Let \(g_{H}\) be the element of \(\mathcal{O}(K)\) defined by \(g_{\mathrm{H}}(z)=z\) in the neighborhood of H and \(g_{\mathrm{H}}(z)=0\) in the neighborhood of K---H. We have \(g_{\mathrm{H}}(z)=f_{\mathrm{H}}(z)z\) on K, and therefore \(x_{\mathrm{H}}=g_{\mathrm{H}}(x)\). It follows that, if \(H\neq K\) , one has

\[
\mathrm{Sp} _ {\mathrm{A}} (x _ {\mathrm{H}}) = g _ {\mathrm{H}} (\mathrm{K}) = \mathrm{H} \cup \{0 \}.
\]

Let \(\lambda \in \mathbf{C} - \mathrm{H}\) . Let \(h_{\mathrm{H},\lambda}\) be the element of \(\mathcal{O}(\mathrm{K})\) equal to \((\lambda - z)^{-1}\) in a neighborhood of H and to 0 in a neighborhood of K-H. We have \(h_{\mathrm{H},\lambda} = f_{\mathrm{H}}h_{\mathrm{H},\lambda}\) and \((\lambda f_{\mathrm{H}} - g_{\mathrm{H}})h_{\mathrm{H},\lambda} = f_{\mathrm{H}}\) . If we denote \(\mathrm{R_H}(x,\lambda) = h_{\mathrm{H},\lambda}(x)\) , we therefore have \(\mathrm{R_H}(x,\lambda) \in \mathrm{B_H}\) and

\[
\begin{array}{r} \mathrm {R_ {H}} (x, \lambda) (\lambda j _ {\mathrm{H}} - x _ {\mathrm{H}}) = (\lambda j _ {\mathrm{H}} - x _ {\mathrm{H}}) \mathrm {R_ {H}} (x, \lambda) = j _ {\mathrm{H}}, \\ \mathrm {R_ {H}} (x, \lambda) j _ {\mathrm {K-H}} = j _ {\mathrm {K-H}} \mathrm {R_ {H}} (x, \lambda) = 0. \end{array}\tag{17}
\]

In particular, \(\lambda \in \mathbf{C} - \mathrm{Sp}_{\mathrm{A_H}}(x_\mathrm{H})\) .

Now let \(\lambda\in\mathrm{H}\) . Suppose that \(\lambda j_{\mathrm{H}}-x_{\mathrm{H}}\) admits an inverse \(y\) in \(\mathrm{A}_{\mathrm{H}}\) . Using the formulas \(j_{\mathrm{H}}y=y\) (since \(y\in\mathrm{A}_{\mathrm{H}}\) ) and \(j_{\mathrm{K-H}}\mathrm{R}_{\mathrm{K-H}}(x,\lambda)=\mathrm{R}_{\mathrm{K-H}}(x,\lambda)\) (since \(\mathrm{R}_{\mathrm{K-H}}(x,\lambda)\in\mathrm{A}_{\mathrm{K-H}}\) ), we find

\[
(\lambda - x) (y + \mathrm{R} _ {\mathrm{K-H}} (x, \lambda)) = (\lambda - x) (j _ {\mathrm{H}} y + j _ {\mathrm{K-H}} \mathrm{R} _ {\mathrm{K-H}} (x, \lambda)) =
\]

\[
(\lambda j _ {\mathrm{H}} - x j _ {\mathrm{H}}) y + (\lambda j _ {\mathrm {K-H}} - x j _ {\mathrm {K-H}}) \mathrm{R} _ {\mathrm {K-H}} (x, \lambda) = j _ {\mathrm{H}} + j _ {\mathrm {K-H}} = 1,
\]

(thanks to formula (17) applied to K---H). One checks similarly that

\[
(y + \mathrm{R} _ {\mathrm{K-H}} (x, \lambda)) (\lambda - x) = 1.
\]

This proves that \(\lambda - x\) admits in A the inverse \(y + \mathrm{R}_{\mathrm{K - H}}(x, \lambda)\) , which is absurd. Thus we have \(\lambda \in \mathrm{Sp}_{\mathrm{A_H}}(x_{\mathrm{H}})\) . We therefore conclude that

\[
\mathrm{Sp} _ {\mathrm{A} _ {\mathrm{H}}} (x _ {\mathrm{H}}) = \mathrm{H}.\tag{18}
\]

In particular, if H is nonempty, the idempotent \(j_{H}\) is nonzero.

Formulas (17) and (18) prove that the function \(\lambda \mapsto \mathrm{R}_{\mathrm{H}}(x,\lambda)\) , defined in \(\mathbf{C} - \mathbf{H}\) , is the resolvent of \(x_{\mathrm{H}}\) with respect to \(\mathrm{A_H}\) .

PROPOSITION 16. --- We keep the previous notation. Let \((\mathrm{H}_{i})_{1\leqslant i\leqslant n}\) be a partition of \(\mathrm{Sp}_{\mathrm{A}}(x)\) into elements of \(\Pi\) .

\begin{enumerate}
\def\labelenumi{\alph{enumi})}
\item
  The algebra B is canonically identified with the algebra \(\mathrm{B_{H_1}}\times \dots \times \mathrm{B_{H_n}}\)
\item
  We have \(x_{\mathrm{H}_i}x_{\mathrm{H}_j} = 0\) for \(i\neq j\) , and
\end{enumerate}

\[
x = x _ {\mathrm{H} _ {1}} + x _ {\mathrm{H} _ {2}} + \dots + x _ {\mathrm{H} _ {n}};
\]

\begin{enumerate}
\def\labelenumi{\alph{enumi})}
\setcounter{enumi}{2}
\tightlist
\item
  We have
\end{enumerate}

\[
\mathrm{R} (x, \lambda) = \mathrm{R} _ {\mathrm{H} _ {1}} (x, \lambda) + \dots + \mathrm{R} _ {\mathrm{H} _ {n}} (x, \lambda)\tag{19}
\]

for all \(\lambda\in\mathbf{C}-\mathrm{Sp}_{\mathrm{A}}(x)\) . In particular, if \(\mathrm{H}\in\Pi\) , the resolvent \(\lambda\mapsto\mathrm{R}(x,\lambda)\) is equal in the neighborhood of \(\mathrm{H}\) to the sum of \(\mathrm{R}_{\mathrm{H}}(x,\lambda)\) and of a holomorphic function.

The relation \(1 = j_{H_{1}} + \cdots + j_{H_{n}}\) is a decomposition of 1 into pairwise orthogonal idempotents of B, so the algebra B is canonically identified with the product algebra \(B_{H_{1}} \times \cdots \times B_{H_{n}}\) (A, I, p.~105, prop. 10).

Assertion \(b\) ) follows from the corresponding relations for the functions \(g_{\mathrm{H}_i}\) ; assertion \(c\) ) is a consequence of \(a\) ) and of the equality \(\mathrm{R}(x_{\mathrm{H}},\lambda) = \mathrm{R}_{\mathrm{H}}(x,\lambda)\) .

PROPOSITION 17. --- Let \(\mu\) be an isolated point of \(\mathrm{Sp}_{\mathrm{A}}(x)\) . Then

\begin{enumerate}
\def\labelenumi{\alph{enumi})}
\tightlist
\item
  For every \(\lambda \in \mathbf{C} - \mathrm{Sp}_{\mathrm{A}}(x)\) , we have
\end{enumerate}

\[
\mathrm{R} (x, \lambda) = \mathrm{R} _ {\{\mu \}} (x, \lambda) + \mathrm{R} _ {\mathrm{Sp} _ {\mathrm{A}} (x) - \{\mu \}} (x, \lambda);
\]

\begin{enumerate}
\def\labelenumi{\alph{enumi})}
\setcounter{enumi}{1}
\item
  The function \(\lambda \mapsto \mathrm{R}_{\mathrm{Sp}_{\mathrm{A}}(x)-\{\mu\}}(x,\lambda)\) is holomorphic in \(\mathbf{C}-\mathrm{Sp}_{\mathrm{A}}(x)\) and in a neighborhood of \(\mu\); moreover, the function \(\lambda \mapsto \mathrm{R}_{\{\mu\}}(x,\lambda)\) is holomorphic in \(\mathbf{C}-\{\mu\}\) ;
\item
  We have
\end{enumerate}

\[
\lim _ {n \to + \infty} \| (x - \mu) ^ {n} j _ {\{\mu \}} \| ^ {1 / n} = 0
\]

and, for \(\lambda \in \mathbf{C} - \{\mu\}\) , the formula

\[
\mathrm{R} _ {\{\mu \}} (x, \lambda) = \sum_ {n = 0} ^ {\infty} (\lambda - \mu) ^ {- n - 1} (x - \mu) ^ {n} j _ {\{\mu \}}.\tag{20}
\]

The preceding implies the assertions \(a\) ) and \(b\) ). Let us prove \(c\) ). By replacing \(x\) by \(x-\mu\) , one reduces to the case where \(\mu=0\) . Set \(H = \{0\}\); this is an open and closed subset of \(\mathrm{Sp}_{\mathrm{A}}(x)\) . By formula (18), the spectrum of \(x_{H}\) in \(\mathrm{A}_{H}\) is \{0\}, hence \(x_{H}\) is quasi-nilpotent, that is, \(\|x^{n}j_{H}\|^{1/n}=\|(xj_{H})^{n}\|^{1/n}\) tends to 0 as \(n\) tends to + \(\infty\) . Moreover, for \(\lambda\neq0\) , in A one has

\[
\left(\lambda j _ {\mathrm{H}} - x _ {\mathrm{H}}\right) ^ {- 1} = \sum_ {n = 0} ^ {\infty} \lambda^ {- n - 1} x _ {\mathrm{H}} ^ {n}
\]

(Theorem 1 of I, p.~24, d)), whence (20).

COROLLARY 1. --- Let \(\mu\) be an isolated point of \(\mathrm{Sp}_{\mathrm{A}}(x)\) and \(p\) be a strictly positive integer. For \(\mu\) to be a pole of order \(p\) of the resolvent of \(x\) (cf.~VAR, R1, p.~30, 3.3.9), it is necessary and sufficient that \((x - \mu)^{p-1} j_{\{\mu\}} \neq 0\) and \((x - \mu)^{p} j_{\{\mu\}} = 0\) .

COROLLARY 2. --- Let \(\mu\) be an isolated point of \(\mathrm{Sp}_{\mathrm{A}}(x)\) . Let \(\Gamma\) be the oriented boundary of an open disk \(\Delta\) with center \(\mu\) such that

\[
\operatorname{Sp} _ {\mathrm{A}} (x) \cap (\Gamma \cup \Delta) = \{\mu \}.
\]

Then the idempotent \(j_{\{\mu \}}\) associated with \(x\) and \(\{\mu\}\) is given by

\[
j _ {\{\mu \}} = \frac {1}{2 i \pi} \int_ {\Gamma} (z - x) ^ {- 1} d z.
\]

*In other words, the idempotent \(j_{\{\mu\}}\) is the residue at \(\mu\) of the resolvent of \(x.*\)

For \(z\in \mathbf{C} - \mathrm{Sp}_{\mathrm{A}}(x)\) , one has

\[
(z - x) ^ {- 1} = \mathrm{R} (x, z) = \mathrm{R} _ {\{\mu \}} (x, z) + \mathrm{R} _ {\mathrm{H}} (x, z),
\]

where H = SpA(x) - \{μ\} (formula (19)). The function z ↦ R\_H(x, z) is holomorphic in C-H and in a neighborhood of \{μ\} (prop. 17, b)), hence

\[
\frac {1}{2 i \pi} \int_ {\Gamma} \mathrm{R} _ {\mathrm{H}} (x, z) d z = 0
\]

(VAR, R2, p.~48, 11.2.5). The function \(z \mapsto \mathrm{R}_{\{\mu\}}(x,z)\) is the resolvent of the element \(j_{\{\mu\}}xj_{\{\mu\}}\) of the unital algebra \(\mathrm{A}_{\{\mu\}}\) . We then have

\[
j _ {\{\mu \}} = \frac {1}{2 i \pi} \int_ {\Gamma} \mathrm{R} _ {\{\mu \}} (x, z) d z
\]

by prop. 9 of I, p.~75 applied to \(A_{\{\mu\}}\) and to the constant function 1 in a neighborhood of \(\Delta \cup \Gamma\) . The corollary follows.

